% command to make sure colorboxes don't create whitespace between code lines
\fboxsep=.6pt

% default settings for listings
\lstset{style=python}

\title[Programmieren]{Programmieren: Verzweigungen}

\subtitle{Negation mit \lstinline|not|}

\author[Wendl, Cyril] % (optional)
{Cyril Wendl}

\institute[KLW] % (optional)
{
    Fachschaft Informatik\\
    Kantonsschule im Lee
}

\date[\today] % (optional)

\logo{\includegraphics[width=2cm]{\thelogo}}

\titlegraphic{\vspace{-.7cm}\includegraphics[width=\textwidth]{"Grundlagen_Info/06_Datenbanken/Figures/title-emojis"}}

\frame{\titlepage}


\begin{frame}[fragile]{Negation von Aussagen}
    \begin{table}
        \begin{tabular}{p{5cm}|p{5cm}}
            \textbf{Aussage} & \textbf{Negation} \\\midrule\midrule
            \visible<+->{``Das Mädchen heisst Elin''}& \visible<+->{``Das Mädchen heisst nicht Elin''}\\\midrule
            \visible<+->{``Niemand in dieser Klasse ist volljährig''}& \visible<+->{``Mindestens eine Person in dieser Klasse ist volljährig''}\\\midrule
            \visible<+->{$x > 2$} & \visible<+->{$x \leq 2$}
        \end{tabular}
    \end{table}
\end{frame}


\begin{frame}[fragile]{Negation von Aussagen in Python: \lstinline|not|}
    \begin{columns}
        \begin{column}{.25\textwidth}
            \begin{lstlisting}
                if x == 0:
                    ...
            \end{lstlisting}
        \end{column}
        \begin{column}{.15\textwidth}
            $\underrightarrow{\text{Negation:}}$
        \end{column}
        \begin{column}{.5\textwidth}
            Entweder:
            \begin{lstlisting}
                if x != 0:
                    ...
            \end{lstlisting}
            
            Oder:
            \begin{lstlisting}
                if not x == 0:
                    ...
            \end{lstlisting}
        \end{column}
    \end{columns}
\end{frame}

\begin{frame}[fragile]{Auftrag}
    \begin{itemize}
        \item \aufgabebox{Aufgaben 5.21 - 5.23}
    \end{itemize}
\end{frame}

\begin{frame}[fragile]{Zusammenfassung}
    \begin{itemize}
        \item Mit einem \lstinline|not| kann ein logischer Ausdruck \textit{negiert} werden
        \item Mit den Befehlen \lstinline|and| sowie \lstinline|or| können logische Ausdrücke miteinander \textit{verknüpft} werden.
    \end{itemize}
\end{frame}



\appendix
\begin{frame}[plain]
    \centering
    \Huge Appendix
\end{frame}


\begin{frame}[fragile]{Mehrere Bedingungen zusammenfassen}{Negation mehrere Bedingunge}
    \begin{lstlisting}[escapechar=!, basicstyle=\ttfamily\footnotesize]
        antwort = input("Nenne einen italienischsprachigen Kanton")
        if !\colorbox{CarnationPink}{not (}!antwort == "Tessin" or antwort == "Graubünden"!\colorbox{CarnationPink}{)}!:
            print("Falsch")
        else:
            print("Richtig")
    \end{lstlisting}
    
    oder:
    
    \begin{lstlisting}[escapechar=!, basicstyle=\ttfamily\footnotesize]
        antwort = input("Nenne einen italienischsprachigen Kanton")
        
        if not antwort == "Tessin" !\colorbox{CarnationPink}{and}! not antwort == "Graubünden":
        print("Falsch")
        else:
        print("Richtig")
    \end{lstlisting}
\end{frame}
